\begin{UPSTIexercice}{Trimestre correspondant à un mois}
    On souhaite écrire un programme qui affiche le trimestre correspondant à un mois donné. Voici le cahier des charges :
    \begin{itemize}
        \item Le mois est stocké dans une variable \texttt{int mois} (compris normalement entre \texttt{1} et \texttt{12}).
        \item Si \texttt{mois} vaut \texttt{1}, \texttt{2} ou \texttt{3}, afficher \texttt{Trimestre 1}.
        \item Si \texttt{mois} vaut \texttt{4}, \texttt{5} ou \texttt{6}, afficher \texttt{Trimestre 2}.
        \item Si \texttt{mois} vaut \texttt{7}, \texttt{8} ou \texttt{9}, afficher \texttt{Trimestre 3}.
        \item Si \texttt{mois} vaut \texttt{10}, \texttt{11} ou \texttt{12}, afficher \texttt{Trimestre 4}.
        \item Pour toute autre valeur de \texttt{mois}, afficher \texttt{Mois invalide}.
        \item \textbf{Contrainte de rédaction :} chacun des cinq messages ci-dessus (\texttt{printf}) ne doit apparaître \textbf{qu'une seule fois} dans le code du programme.
    \end{itemize}
    \UPSTIquestion{Écrire, en langage C, la structure \texttt{switch} correspondant à ce cahier des charges, en respectant la contrainte de rédaction. On testera le programme avec \texttt{mois = 8}.}
\end{UPSTIexercice}

\begin{UPSTIprofOnlyEnv}
    \begin{UPSTIcorrectionP}{Trimestre correspondant à un mois}
        \begin{lstlisting}[language=c]
int mois = 8;

switch (mois)
{
    case 1:
    case 2:
    case 3:
        printf("Trimestre 1\n");
        break;
    case 4:
    case 5:
    case 6:
        printf("Trimestre 2\n");
        break;
    case 7:
    case 8:
    case 9:
        printf("Trimestre 3\n");
        break;
    case 10:
    case 11:
    case 12:
        printf("Trimestre 4\n");
        break;
    default:
        printf("Mois invalide\n");
        break;
}
        \end{lstlisting}
        Avec \texttt{mois = 8}, l'affichage obtenu est~: \texttt{Trimestre 3}.

        Point de rigueur essentiel~: la contrainte \og chaque \texttt{printf} qu'une seule fois \fg{} interdit d'écrire, par exemple, \texttt{case 1:}, \texttt{case 2:} et \texttt{case 3:} chacun suivi de sa propre copie de \texttt{printf("Trimestre 1\textbackslash n"); break;}. La seule façon de regrouper trois \texttt{case} sur un même bloc, en C, est le \textbf{fall-through} : plusieurs \texttt{case} \textbf{vides} (sans instruction entre eux) qui \og tombent \fg{} ensemble dans le même bloc d'instructions. L'énoncé ne nommait pas cette technique, mais la contrainte de rédaction la rendait obligatoire.

        Non, la contrainte ne visait pas spécifiquement \texttt{switch} : une cascade \texttt{if}/\texttt{else if} peut tout aussi bien respecter la règle \og un seul \texttt{printf} par message \fg{}, en combinant les conditions avec \texttt{||} (par exemple \texttt{if (mois == 1 || mois == 2 || mois == 3))}. Le fall-through est \textbf{l'équivalent}, côté \texttt{switch}, du \texttt{||} côté \texttt{if} : les deux techniques répondent à la même contrainte, chacune dans la syntaxe de sa propre structure.
    \end{UPSTIcorrectionP}
\end{UPSTIprofOnlyEnv}
